\section{Introduction}

\subsection{DALI Overview}
DALI (Dynamic Agent Logic framework for Interaction) is a logic-based multi-agent system
built on top of SICStus Prolog. Each agent is defined by a set of \emph{external events}
(reactions to incoming messages), \emph{internal events} (proactive behaviours triggered by
conditions), and standard Prolog clauses for knowledge representation and reasoning.

DALI agents run as independent SICStus processes that communicate via a shared message
server (\texttt{LINDA}) running on a configurable TCP port.

\subsection{Project Goals}
The F1 Race simulation was developed with the following goals in mind:
\begin{itemize}[leftmargin=2em]
  \item Demonstrate the \emph{reactive} and \emph{proactive} capabilities of DALI agents in
        a realistic, event-driven scenario.
  \item Show how a fully \emph{dynamic} agent roster can be maintained through code
        generation, avoiding any hardcoded agent names in the DALI source.
  \item Provide a real-time monitoring interface that reflects agent activity without
        modifying the agent logic.
  \item Package the entire stack as a reproducible Docker Compose application.
\end{itemize}
